\documentclass[10pt,a4paper]{amsart}
\usepackage[margin=19mm]{geometry}
\usepackage{amsmath,amssymb,mathtools}
\usepackage[scr=rsfs,cal=euler]{mathalfa}
\usepackage{xfrac}
\usepackage[unicode,hidelinks,bookmarks=false]{hyperref}
\newcommand{\ie}{i.e.\ }
\newcommand{\cf}{cf.\ }
\newcommand{\resp}{respectively\ }
\newcommand{\Subsection}{\subsection}
\providecommand{\nobreakdash}{\nobreak\hbox{-}}
\setlength{\emergencystretch}{2em}
\multlinegap0pt
\usepackage{amssymb}
\usepackage{mathtools}
\usepackage[all]{xy}
\usepackage{tikz}
\newtheorem{proposition}{Proposition}
\newcommand{\e}{\epsilon}
\title{Quantum $K$-invariants and Gopakumar--Vafa invariants II. \\ Calabi-Yau threefolds at genus zero}
\author{Y.-C.~Chou, Y.-P.~Lee}
\date{Source: arXiv:2305.08480v4 (CC BY 4.0)}
\begin{document}
\maketitle

\noindent\emph{Source and licence.} Quantum $K$-invariants and Gopakumar--Vafa invariants II. \\ Calabi-Yau threefolds at genus zero. Version arXiv:2305.08480v4 (CC BY 4.0), distributed by arXiv under the Creative Commons Attribution 4.0 International licence, http://creativecommons.org/licenses/by/4.0/, as recorded in the arXiv licence metadata for this exact version and shown on the abstract page https://arxiv.org/abs/2305.08480v4. Full licence notice: \texttt{licenses/arxiv-2305.08480-CC-BY-4.0.txt}.\par\smallskip

\noindent\textit{Editorial scope.} Counted unit: the article's proposition prop\_multicover\_t (source0.tex lines 1456--1468). The article's complete original proof of it is delivered with it (source0.tex lines 1469--1507, one \texttt{proof} environment of 39 lines). No mathematics is written, repaired or re-derived: the statement and the proof are the article's own lines, in the article's own notation, and no retained line is substituted or re-flowed.  No further source-local context is needed: every cross-reference of every retained line resolves inside the two delivered spans.  Nothing was omitted from any retained span, so the package carries no omission record.  No retained line contains a citation, so the wrapper carries no bibliography and no bibliography entry.  No retained line contains a figure environment, an image-inclusion command or drawing code, so no image is delivered and the package declares no asset dependency.  Editorial formatting only, in addition: an AMS \texttt{amsart} preamble that restates the article's own declarations and macros used by the retained lines, namely the environments \texttt{proposition} on separate counters, together with the standard packages those lines need.  The retained environments print in this document as Proposition 1 in order of appearance, whereas the article prints the same items with the numbering of its own full document.  The complete native source of the article as distributed in the arXiv e-print for this exact version is bundled unchanged as \texttt{sources/141493/source0.tex}.
\begin{proposition} \label{prop_multicover_t}
The $J$-function for %(any toric completion of)  
$Y_{-1,-1}$, with curve classes in the zero section, is 
\begin{equation} \label{e:4.3}
\begin{split}
\frac{1}{1-q} & J^{Y_{-1,-1}}(t, q, Q) = 1 + \frac{t_1(1-P)}{1-q}
\\
&+  (1-PT)^2\Big(1 +(1-P)\Big)\sum_{r \geq 1} \, Q^r\Big( a(r, q^r) + c(r,q, t_1) \Big) % Q^r \Big( \frac{r-1}{1-q^r} + \frac{1}{(1-q^r)^2} \Big)
\\
& + (1-PT)^2(1-P) \sum_{r\geq 1} Q^r \, \Big( b(r, q^r) + d(r,q, t_1) \Big) . %\Big( \frac{r^2-1}{1-q^r} + \frac{3}{(1-q^r)^2} -\frac{2}{(1-q^r)^3} \Big) .
\end{split} 
\end{equation}
\end{proposition}

\begin{proof}
We compute $J^{Y_{-1,-1}}$ in two steps. Consider the flow:
\[
\begin{split}
    J^{Y_{-1,-1}}(t^*,q,Q) =\sum_{r\geq 0} Q^r \exp \left( t_1 \frac{1-Pq^r}{1-q} \right) J^{Y_{-1,-1}}(0,q,Q),
\end{split}
\]
where 
\[
    t^* = t + (1-PT)^2(E_1(t_1,q,Q)) + (1-P)(1-PT)^2(E_2(t_1,q,Q))
\]
with
\[
\begin{split}
    E_1(t_1,q,Q) & \in F_1(t_1,Q) + (1-q) \cdot  \mathbb{Q}[t_1,q^{-1},q][\![Q]\!];
    \\
    E_2(t_1,q,Q) & \in \mathbb{Q}[t_1,q^{-1},q][\![Q]\!].
\end{split}
\]
Then we consider the explicit reconstruction
\[
\begin{split}
    &J^{Y_{-1,-1}}(t, q, Q) =  \sum_{r\geq 0}  J_r^{Y_{-1,-1}}(t^*,q,Q) Q^r  
    \\ 
    & \cdot\exp \left( \frac{ \sum_{i=0}^1\delta_i(t_1,Q)(1-Pq^{r})(1-PT q^r)^i + \sum_{i=0}^3 \e_i(t_1,Q) (1-PT q^{r})^i}{(1-q)} \right) 
    \\
    &\quad \cdot\left( \sum_{i=0}^1 s_i(t_1,q,Q)(1-Pq^r)(1-PT q^r)^i + \sum_{i=0}^3 u_i(t_1,q,Q) (1-PTq^r)^i \right),
\end{split}
\]
for some uniquely determined $\delta_i(t_1,Q)$, $\epsilon_i(t_1,Q)$, $s_i(t_1,q,Q)$, and $u_i(t_1,q,Q)$, where
\[
\begin{split}
    \epsilon_2(t_1,Q) &= - F_1(t_1,Q), \quad u_0(t_1,q,Q) = 1,
    \\
    u_1(t_1,q,Q) &=  \epsilon_i(t_1,Q)=\delta_i(t_1,Q)=s_i(t_1,q,Q)=0\quad \forall i=0,1.
\end{split}
\]
Note that for any choice of $\epsilon_3(t_1,Q)$, $u_2(t_1,q,Q)$, and $u_3(t_1,q,Q)$ will not change the $\mathcal{K}_-$ part. The proposition follows from a long and direct computation.
\end{proof}

\end{document}
